Com $F_0=10\,000$, $E_0=5\,000$ e $\lambda_E=\lambda_F=\lambda=0{,}1$ por dia, o parâmetro $\alpha = \sqrt{\lambda^2} = 0{,}1$.

\textbf{Vencedor:} O exército com $F_0 > E_0$ vence, pois pela lei de Lanchester quadrática, a condição de vitória do exercito $F$ é $\lambda_F F_0^2 > \lambda_E E_0^2$, ou seja, $10\,000^2 > 5\,000^2$. \textbf{$F$ vence.}

\textbf{Tempo de batalha:}
Da solução $E(t)=0$:
\[
\tanh(\alpha t_c) = \sqrt{\frac{\lambda_F}{\lambda_E}}\frac{F_0}{E_0} = 1\cdot\frac{10\,000}{5\,000} = 2.
\]
Porém $|\tanh(x)|<1$ para todo $x$ finito. Portanto, a solução exata para um
$\tanh = 2 > 1$ não existe em tempo finito com $\tanh$ puro: usa-se a fórmula alternativa da condição de vitória e a solução do Problema 5.31 dando:
\[
t_c = \frac{1}{\alpha}\operatorname{arctanh}\!\left(\frac{E_0/F_0}{\sqrt{\lambda_F/\lambda_E}}\right)=\frac{1}{0{,}1}\operatorname{arctanh}(0{,}5)\approx 5{,}49\ \mathrm{dias}.
\]

\textbf{Tropas restantes de $F$:}
\[
F(t_c) = \sqrt{F_0^2 - \frac{\lambda_E}{\lambda_F}E_0^2} = \sqrt{10\,000^2 - 5\,000^2} = \sqrt{75\,000\,000} \approx 8\,660\ \text{soldados}.
\]